\documentclass[aspectratio=169,10pt]{beamer}
% Shared preamble for the Algorithms and Programming Concepts slides (Beamer, Metropolis)
\usetheme[progressbar=frametitle,numbering=fraction,background=light,block=fill,sectionpage=none]{metropolis}
\usepackage[T1]{fontenc}
\usepackage{FiraSans}
\usepackage[scaled=.85]{FiraMono}
\usepackage{amsmath,amssymb}
\usefonttheme[onlymath]{serif}
\usepackage{booktabs,array,tabularx}
\usepackage{tikz,pgfplots}
\usepackage{xcolor}
\usepackage{listings}
\usepackage{adjustbox}
\usepackage{fontawesome5}
\hyphenpenalty=10000\exhyphenpenalty=10000
\newenvironment{sfig}[1][0.80]{\begin{adjustbox}{max width=\linewidth,max totalheight=#1\textheight,center}}{\end{adjustbox}}
\usetikzlibrary{arrows.meta,positioning,shapes.geometric,shapes.misc,fit,calc,backgrounds,decorations.pathreplacing,patterns,matrix,chains}
\pgfplotsset{compat=1.17}

\definecolor{tdInk}{HTML}{1D2330}
\definecolor{tdMuted}{HTML}{5D6678}
\definecolor{tdLine}{HTML}{D9D5CA}
\definecolor{tdPaper}{HTML}{FBFAF7}
\definecolor{tdBlue}{HTML}{2457D6}
\definecolor{tdGreen}{HTML}{17855A}
\definecolor{tdAmber}{HTML}{B7791F}
\definecolor{tdRed}{HTML}{C8362F}
\definecolor{tdPurple}{HTML}{7A4FD1}
\setbeamercolor{normal text}{fg=tdInk,bg=tdPaper}
\setbeamercolor{alerted text}{fg=tdRed}
\setbeamercolor{example text}{fg=tdGreen}
\setbeamercolor{frametitle}{bg=tdInk,fg=white}
\setbeamercolor{progress bar}{fg=tdBlue,bg=tdLine}
\setbeamercolor{title separator}{fg=tdBlue}
\setbeamercolor{block title}{bg=tdBlue!12,fg=tdInk}
\setbeamercolor{block body}{bg=tdBlue!5,fg=tdInk}
\setbeamercolor{block title alerted}{bg=tdRed!12,fg=tdRed!80!black}
\setbeamercolor{block body alerted}{bg=tdRed!5,fg=tdInk}
\setbeamercolor{block title example}{bg=tdGreen!14,fg=tdGreen!70!black}
\setbeamercolor{block body example}{bg=tdGreen!6,fg=tdInk}
\setbeamerfont{frametitle}{size=\large,series=\bfseries}
\makeatletter\setlength{\metropolis@frametitle@padding}{1.15ex}\makeatother
\setbeamertemplate{navigation symbols}{}
\setbeamerfont{block title}{size=\small}

% ---- code listings (Python 3.14) ----
\lstdefinestyle{py}{language=Python,
  basicstyle=\ttfamily\fontsize{8.4}{10.1}\selectfont\color{tdInk},
  keywordstyle=\color{tdBlue}\bfseries, commentstyle=\color{tdMuted}\itshape, stringstyle=\color{tdGreen},
  morekeywords={match,case,type,self,True,False,None},
  emph={print,len,range,sorted,min,max,sum,enumerate,zip,heappush,heappop,cache},emphstyle=\color{tdPurple},
  showstringspaces=false, columns=fullflexible, keepspaces=true, frame=single, rulecolor=\color{tdLine},
  backgroundcolor=\color{white}, xleftmargin=2pt, xrightmargin=2pt, aboveskip=2pt, belowskip=0pt, framesep=3pt}
\lstdefinestyle{out}{basicstyle=\ttfamily\fontsize{8}{9.6}\selectfont\color{tdInk}, columns=fullflexible, keepspaces=true,
  frame=single, rulecolor=\color{tdGreen!60}, backgroundcolor=\color{tdGreen!6}, xleftmargin=2pt, xrightmargin=2pt, aboveskip=2pt, belowskip=0pt, framesep=3pt}
% \pycode[size]{snippet} shows snippets/<snippet>.py ; \pyout[size]{snippet} shows snippets/<snippet>.out
\newcommand{\pycode}[2][\relax]{{#1\lstinputlisting[style=py]{snippets/#2.py}}}
\newcommand{\pyout}[2][\relax]{{#1\lstinputlisting[style=out]{snippets/#2.out}}}

% ---- tikz styles ----
\tikzset{>={Stealth[length=2.2mm]},
  card/.style={rounded corners=4pt,draw=#1!70!black,fill=#1!10,align=center,inner sep=5pt,font=\small}, card/.default=tdBlue,
  pill/.style={rounded corners=7pt,draw=none,fill=#1,text=white,inner sep=3pt,font=\scriptsize\bfseries,align=center}, pill/.default=tdBlue,
  arr/.style={->,thick,draw=tdMuted},
  cell/.style={draw=tdInk!70,minimum width=7.5mm,minimum height=7.5mm,inner sep=0pt,font=\small\ttfamily,fill=white},
  tn/.style={circle,draw=tdInk!80,fill=white,minimum size=7mm,inner sep=0pt,font=\small\bfseries}}
\newcommand{\key}[1]{\textbf{\color{tdBlue}#1}}
\newcommand{\bad}[1]{\textbf{\color{tdRed}#1}}
\newcommand{\good}[1]{\textbf{\color{tdGreen!70!black}#1}}
\newcommand{\apcfooter}[1]{\setbeamertemplate{footline}{\hfill\tiny\color{tdMuted}Algorithms and Programming Concepts \,|\, #1 \,|\, BUV\hspace{1.6em}\insertframenumber/\inserttotalframenumber\hspace{0.8em}\vspace{4pt}}}
\newcommand{\apctitle}[2]{\title{#1}\subtitle{#2}\author{Dr.\ Quang-Vinh Dang\\ {\small Senior Lecturer (Assoc.\ Prof)}}\institute{British University Vietnam}\date{Algorithms and Programming Concepts \,|\, Year 1 \,|\, Python 3.14}}

\apcfooter{Trees}
\apctitle{Trees}{Hierarchies, binary search trees, traversals and heaps}

% seven-node binary tree; fill colours \fA..\fG (set before use), labels as arguments
\newcommand{\fA}{white}\newcommand{\fB}{white}\newcommand{\fC}{white}\newcommand{\fD}{white}
\newcommand{\fE}{white}\newcommand{\fF}{white}\newcommand{\fG}{white}
\newcommand{\bsttree}[7]{%
  \node[tn,fill=\fA] (a) at (0,0) {#1};
  \node[tn,fill=\fB] (b) at (-2,-1.15) {#2};  \node[tn,fill=\fC] (c) at (2,-1.15) {#3};
  \node[tn,fill=\fD] (d) at (-3,-2.3) {#4};   \node[tn,fill=\fE] (e) at (-1,-2.3) {#5};
  \node[tn,fill=\fF] (f) at (1,-2.3) {#6};    \node[tn,fill=\fG] (g) at (3,-2.3) {#7};
  \foreach \p/\q in {a/b,a/c,b/d,b/e,c/f,c/g}{\draw[tdInk!60] (\p) -- (\q);}}
\newcommand{\drawnodes}[7]{\foreach \nm/\lab in {a/#1,b/#2,c/#3,d/#4,e/#5,f/#6,g/#7}{\node[tn,fill=white] at (\nm) {\lab};}}

\begin{document}
\begin{frame}[plain]\titlepage\end{frame}

\begin{frame}{Trees are everywhere}
\begin{columns}[T]
\begin{column}{0.5\textwidth}
{\small
Lists are one-dimensional. Many things in computing are \key{hierarchies}:
\begin{itemize}\setlength\itemsep{3pt}
\item the \textbf{folders} and files on your disk;
\item the \textbf{HTML} structure of a web page (the DOM);
\item an \textbf{organisation chart} or a family tree;
\item a \textbf{decision tree} in machine learning;
\item the \textbf{expression} $(3+4)\times(10-2)$ as a compiler sees it;
\item a \textbf{priority queue} that always gives the most urgent job.
\end{itemize}}
\end{column}
\begin{column}{0.46\textwidth}
\begin{block}{After this lecture you can}
{\small
\begin{itemize}\setlength\itemsep{2pt}
\item use the vocabulary: root, leaf, height, subtree;
\item build and search a \textbf{binary search tree};
\item walk a tree in \textbf{pre-, in-, post- and level order};
\item explain why \textbf{height} decides the speed;
\item use \texttt{heapq} as a priority queue.
\end{itemize}}
\end{block}
\end{column}
\end{columns}
\end{frame}

\begin{frame}{Vocabulary}
\begin{columns}[T]
\begin{column}{0.5\textwidth}
\begin{sfig}[0.7]
\begin{tikzpicture}
  \begin{scope}[on background layer]
    \fill[tdBlue!10,rounded corners=8pt] (0.6,-0.5) rectangle (3.9,-3.0);
  \end{scope}
  \bsttree{A}{B}{C}{D}{E}{F}{G}
  \node[font=\scriptsize,text=tdRed,anchor=west] at (0.55,0.45) {root};
  \node[font=\scriptsize,text=tdGreen!70!black,anchor=north] at (-3,-2.75) {leaf};
  \node[font=\scriptsize,text=tdGreen!70!black,anchor=north] at (3,-2.75) {leaf};
  \node[font=\scriptsize,text=tdBlue,anchor=west] at (3.15,-0.7) {subtree rooted at C};
  \foreach \l/\y in {0/0,1/-1.15,2/-2.3}{\node[font=\tiny,text=tdMuted] at (-4.3,\y) {level \l};}
\end{tikzpicture}
\end{sfig}
\end{column}
\begin{column}{0.46\textwidth}
{\small
\begin{description}\setlength\itemsep{2pt}
\item[node] one item in the tree.
\item[edge] a link from a \textbf{parent} to a \textbf{child}: B is the parent of D and E.
\item[root] the one node with no parent (A).
\item[leaf] a node with no children (D, E, F, G).
\item[subtree] a node together with all its descendants.
\item[depth / level] number of edges from the root: A is at level 0.
\item[height] the number of levels (here 3).
\end{description}
In a \key{binary tree} every node has \textbf{at most two} children: left and right.}
\end{column}
\end{columns}
\end{frame}

\begin{frame}{A binary tree in Python}
\begin{columns}[T]
\begin{column}{0.5\textwidth}
\pycode{node}
\vspace{3mm}
{\small
A node stores a \textbf{key} and links to its \textbf{left} and \textbf{right} children. \texttt{None} means ``no child here''. A tree is just its root node: \texttt{Node(50, Node(30), Node(70))}.}
\end{column}
\begin{column}{0.46\textwidth}
\begin{block}{Python 3.14 note}
{\small
Inside \texttt{class Node}, the annotation \texttt{Node | None} refers to the class itself. Before 3.14 you had to write it as a string, \texttt{"Node | None"}, because \texttt{Node} did not exist yet. Python 3.14 evaluates annotations \textbf{lazily}, so the plain form works.}
\end{block}
\vspace{2mm}
{\small
\texttt{@dataclass(slots=True)} writes \texttt{\_\_init\_\_} and \texttt{\_\_repr\_\_} for us and stores the fields compactly, which helps when a tree has millions of nodes.}
\end{column}
\end{columns}
\end{frame}

\begin{frame}{Binary search tree: keeping the keys in order}
\begin{columns}[T]
\begin{column}{0.5\textwidth}
\begin{sfig}[0.62]
\begin{tikzpicture}
  \renewcommand{\fA}{tdAmber!40}\renewcommand{\fC}{tdAmber!40}\renewcommand{\fF}{tdGreen!35}
  \bsttree{50}{30}{70}{20}{40}{60}{80}
  \draw[tdAmber!80!black,very thick,->] (a) -- (c); \draw[tdAmber!80!black,very thick,->] (c) -- (f);
  \node[font=\scriptsize,text=tdAmber!70!black,anchor=west] at (0.55,0.45) {$60>50$: go right};
  \node[font=\scriptsize,text=tdAmber!70!black,anchor=west] at (3.2,-1.1) {$60<70$: go left};
  \node[font=\scriptsize,text=tdGreen!70!black,anchor=north] at (1,-2.7) {found 60};
\end{tikzpicture}
\end{sfig}
\end{column}
\begin{column}{0.48\textwidth}
\begin{block}{The BST rule}
{\small For every node: \textbf{all keys in its left subtree are smaller, all keys in its right subtree are larger}.}
\end{block}
\vspace{2mm}
{\small
Searching for 60: compare with the root (50), go right; compare with 70, go left; found. Each comparison discards a whole \textbf{subtree}, just like binary search discards half of an array.

\vspace{2mm}
\textbf{Inserting} follows the same path and places the new key where the search ends.}
\end{column}
\end{columns}
\end{frame}

\begin{frame}{Insert and search in Python}
\begin{columns}[T]
\begin{column}{0.52\textwidth}
\pycode{bst}
\end{column}
\begin{column}{0.44\textwidth}
\pycode{build}
\vspace{2mm}
\pyout{bst_demo}
\vspace{2mm}
{\small
\texttt{insert} is \textbf{recursive}: it returns the (possibly new) root of the subtree. \texttt{search} uses a loop. Duplicates are ignored. \texttt{build} inserts keys one by one.}
\end{column}
\end{columns}
\end{frame}

\begin{frame}{Four ways to walk a tree}
\begin{sfig}[0.66]
\begin{tikzpicture}[scale=0.8,transform shape]
  \node[font=\bfseries] at (0,0.8) {preorder: node, left, right};
  \bsttree{1}{2}{5}{3}{4}{6}{7}
  \begin{scope}[xshift=9cm]
  \node[font=\bfseries] at (0,0.8) {inorder: left, node, right};
  \bsttree{4}{2}{6}{1}{3}{5}{7}
  \end{scope}
  \begin{scope}[yshift=-4.4cm]
  \node[font=\bfseries] at (0,0.8) {postorder: left, right, node};
  \bsttree{7}{3}{6}{1}{2}{4}{5}
  \end{scope}
  \begin{scope}[xshift=9cm,yshift=-4.4cm]
  \node[font=\bfseries] at (0,0.8) {level order: level by level};
  \bsttree{1}{2}{3}{4}{5}{6}{7}
  \end{scope}
\end{tikzpicture}
\end{sfig}
\vspace{1mm}
{\footnotesize The numbers show the \textbf{order of visit} on the tree with keys $50,30,70,20,40,60,80$. The first three are \emph{depth-first} (recursion); level order is \emph{breadth-first} (a queue).}
\end{frame}

\begin{frame}{Traversals in Python}
\begin{columns}[T]
\begin{column}{0.58\textwidth}
\pycode[\fontsize{7.5}{9}\selectfont]{traversals}
\end{column}
\begin{column}{0.4\textwidth}
\pyout{traversal_demo}
\vspace{2mm}
{\small
\begin{itemize}\setlength\itemsep{2pt}
\item \textbf{inorder} of a BST gives the keys \textbf{sorted}: a free sort!
\item \textbf{preorder} copies a tree.
\item \textbf{postorder} visits children first: deleting a tree, evaluating an expression.
\item \textbf{level order} uses a \texttt{deque} as a queue.
\end{itemize}}
\end{column}
\end{columns}
\end{frame}

\begin{frame}{Height decides the speed}
\begin{sfig}[0.36]
\begin{tikzpicture}
  \renewcommand{\fA}{tdGreen!22}\renewcommand{\fB}{tdGreen!22}\renewcommand{\fC}{tdGreen!22}\renewcommand{\fD}{tdGreen!22}\renewcommand{\fE}{tdGreen!22}\renewcommand{\fF}{tdGreen!22}\renewcommand{\fG}{tdGreen!22}
  \bsttree{50}{30}{70}{20}{40}{60}{80}
  \node[font=\small\bfseries,text=tdGreen!70!black] at (0,0.8) {balanced: height 3};
  \begin{scope}[xshift=9.5cm,yshift=0.2cm]
  \foreach \k [count=\i from 0] in {10,20,30,40,50,60,70}{\node[tn,fill=tdRed!15] (c\i) at (\i*0.9-2.7,-\i*0.62+0.1) {\k};}
  \foreach \i in {1,...,6}{\pgfmathtruncatemacro{\j}{\i-1}\draw[tdInk!60] (c\j) -- (c\i);}
  \node[font=\small\bfseries,text=tdRed] at (0,0.8) {inserted in sorted order: height 7};
  \end{scope}
\end{tikzpicture}
\end{sfig}
\vspace{1mm}
\begin{columns}[T]
\begin{column}{0.47\textwidth}
\pyout{height_demo}
\vspace{1mm}
{\small
Search, insert and delete follow \textbf{one path from the root}, so they cost $O(\text{height})$: balanced $\Rightarrow$ $O(\log n)$, a chain $\Rightarrow$ $O(n)$, no better than a list.}
\end{column}
\begin{column}{0.49\textwidth}
\pyout{height_scale}
\vspace{1mm}
{\small
Random keys give a tree that is not perfect but still shallow. \textbf{Self-balancing trees} (AVL, red--black) rotate nodes after each insert to keep the height $O(\log n)$; databases use a relative, the \textbf{B-tree}.}
\end{column}
\end{columns}
\end{frame}

\begin{frame}{Expression trees and recursion with \texttt{match}}
\begin{columns}[T]
\begin{column}{0.46\textwidth}
\begin{sfig}[0.55]
\begin{tikzpicture}
  \node[tn,fill=tdAmber!30] (r) at (0,0) {$\times$};
  \node[tn,fill=tdAmber!30] (p) at (-1.8,-1.2) {$+$};  \node[tn,fill=tdAmber!30] (m) at (1.8,-1.2) {$-$};
  \node[tn] (a) at (-2.6,-2.4) {3}; \node[tn] (b) at (-1,-2.4) {4}; \node[tn] (c) at (1,-2.4) {10}; \node[tn] (d) at (2.6,-2.4) {2};
  \foreach \u/\v in {r/p,r/m,p/a,p/b,m/c,m/d}{\draw[tdInk!60] (\u) -- (\v);}
  \node[font=\scriptsize,text=tdMuted] at (-2.1,-0.35) {$3+4=7$}; \node[font=\scriptsize,text=tdMuted] at (2.1,-0.35) {$10-2=8$};
  \node[font=\scriptsize,text=tdGreen!70!black] at (0,-3.15) {root: $7\times 8=56$};
\end{tikzpicture}
\end{sfig}
\vspace{1mm}
\pyout{expr_demo}
\end{column}
\begin{column}{0.5\textwidth}
\pycode{expr}
\vspace{2mm}
{\small
Internal nodes are operators, leaves are numbers. To evaluate a node, first evaluate its children, then combine: that is a \textbf{postorder} walk.

\vspace{1mm}
\texttt{match} (Python 3.10+) takes a tree apart by its \emph{shape}: \texttt{("+", a, b)} matches a 3-tuple whose first item is \texttt{"+"} and binds \texttt{a} and \texttt{b}.}
\end{column}
\end{columns}
\end{frame}

\begin{frame}{Heaps and priority queues}
\begin{columns}[T]
\begin{column}{0.46\textwidth}
\begin{sfig}[0.55]
\begin{tikzpicture}
  \node[tn,fill=tdAmber!35] (a) at (0,0) {1};
  \node[tn] (b) at (-1.6,-1.15) {3}; \node[tn] (c) at (1.6,-1.15) {2};
  \node[tn] (d) at (-2.4,-2.3) {5}; \node[tn] (e) at (-0.8,-2.3) {9}; \node[tn] (f) at (0.8,-2.3) {8};
  \foreach \u/\v in {a/b,a/c,b/d,b/e,c/f}{\draw[tdInk!60] (\u) -- (\v);}
  \foreach \v [count=\i from 0] in {1,3,2,5,9,8}{\node[cell,minimum width=6.5mm,fill=white] at (\i*0.68-1.7,-3.5) {\scriptsize\v};}
  \node[font=\tiny,text=tdMuted,anchor=east] at (-2.1,-3.5) {list:};
\end{tikzpicture}
\end{sfig}
{\small A \textbf{min-heap}: every parent is $\le$ its children, so the smallest item is always at the root. It is stored in a plain list.}
\end{column}
\begin{column}{0.5\textwidth}
\pycode{heap}
\vspace{2mm}
\pyout{heap_demo}
\vspace{2mm}
{\small
\texttt{heappush} and \texttt{heappop} cost $O(\log n)$; peeking at \texttt{tasks[0]} is $O(1)$. Use a heap whenever you repeatedly need ``the smallest/most urgent item'': job schedulers, Dijkstra's algorithm, top-$k$.}
\end{column}
\end{columns}
\end{frame}

\begin{frame}{Costs at a glance}
\renewcommand{\arraystretch}{1.32}\footnotesize
\begin{tabular}{@{}>{\raggedright\bfseries}p{42mm}>{\raggedright}p{32mm}>{\raggedright}p{30mm}>{\raggedright\arraybackslash}p{38mm}@{}}
\toprule
Structure & Search & Insert & Notes\\\midrule
Unsorted list & $O(n)$ & $O(1)$ at the end & simple\\
Sorted list (\texttt{bisect}) & $O(\log n)$ & $O(n)$ & shifting items\\
BST, balanced & $O(\log n)$ & $O(\log n)$ & needs rebalancing\\
BST, degenerate chain & $O(n)$ & $O(n)$ & keys inserted in order\\
Heap (\texttt{heapq}) & $O(n)$ any key & $O(\log n)$ & min in $O(1)$\\
Hash table (\texttt{set}, \texttt{dict}) & $O(1)$ average & $O(1)$ average & no order\\
\bottomrule
\end{tabular}
\vspace{4mm}

{\small A tree is a good choice when you need the data \textbf{in order} and also fast insertion and lookup.}
\end{frame}

\begin{frame}{Recap and practice}
\begin{columns}[T]
\begin{column}{0.48\textwidth}
\begin{block}{Take-aways}
{\small
\begin{enumerate}\setlength\itemsep{3pt}
\item A tree is a \key{recursive} structure: a node plus subtrees.
\item In a \key{BST}, left $<$ node $<$ right, so search and insert follow one path.
\item \key{Inorder} traversal of a BST lists the keys in sorted order.
\item Speed depends on the \textbf{height}: $\log n$ if balanced, $n$ if it degenerates.
\item A \key{heap} gives the minimum in $O(1)$ and push/pop in $O(\log n)$.
\end{enumerate}}
\end{block}
\end{column}
\begin{column}{0.48\textwidth}
\begin{block}{Practice}
{\small
\begin{enumerate}\setlength\itemsep{3pt}
\item Draw the BST obtained by inserting \texttt{[8, 3, 10, 1, 6, 14]} and write its four traversals.
\item Write \texttt{minimum(root)} and \texttt{maximum(root)} for a BST.
\item Write \texttt{count\_leaves(root)} with recursion.
\item Add a \texttt{"/"} case to \texttt{evaluate} and print \texttt{("/", 20, ("+", 1, 4))}.
\item Use \texttt{heapq} to get the 3 largest items of a list (hint: \texttt{heapq.nlargest}). How does its speed compare to sorting?
\end{enumerate}}
\end{block}
\end{column}
\end{columns}
\end{frame}
\end{document}
